% !TEX root = ../main.tex

Il file \textit{io\_utils.py} raccoglie le funzioni di caricamento dati. Espone un'unica funzione, \texttt{load\_csv}, che avvolge \texttt{pandas.read\_csv} mantenendo una surface piccola e stabile per notebook e script didattici.

\section{\texorpdfstring{load\_csv(path, *, sep, decimal, rename\_columns, required\_columns, missing, comment, skip\_initial\_space)}{load\_csv(path, *, sep, decimal, rename_columns, required_columns, missing, comment, skip_initial_space)}}
\label{sec:load-csv}
\begin{apidesc}
  \item[\textbf{Scopo}]
    Caricare un file CSV in un \texttt{DataFrame}, con supporto a separatori e decimali personalizzati, rinomina delle colonne, controllo delle colonne richieste e policy esplicita per i valori mancanti.
  \item[\textbf{Parametri}]\hfill
    \begin{description}[font=\normalfont\ttfamily\bfseries, labelwidth=9em, leftmargin=9.5em, itemsep=1pt, parsep=0pt, topsep=2pt]
      \item[path] \texttt{str} o \texttt{Path} $\to$ percorso del file CSV
      \item[sep] \texttt{str}, default \texttt{","} $\to$ separatore di colonna
      \item[decimal] \texttt{str}, default \texttt{"."} $\to$ separatore decimale
      \item[rename\_columns] mapping o \texttt{None} $\to$ mappa \texttt{\{vecchio: nuovo\}} per rinominare le colonne
      \item[required\_columns] collezione o \texttt{None} $\to$ nomi di colonna che devono essere presenti dopo l'eventuale rinomina
      \item[missing] \texttt{"error"} $\mid$ \texttt{"drop"} $\mid$ \texttt{"allow"} (default \texttt{"error"}) $\to$ politica di gestione dei valori mancanti
      \item[comment] \texttt{str} o \texttt{None}, default \texttt{None} $\to$ carattere di commento passato a \texttt{pandas.read\_csv}
      \item[skip\_initial\_space] \texttt{bool}, default \texttt{True} $\to$ se \texttt{True}, ignora gli spazi dopo il separatore
    \end{description}
  \item[\textbf{Note}]
    \texttt{required\_columns} viene verificato dopo \texttt{rename\_columns}. I parametri dopo \texttt{path} sono tutti keyword-only.
  \item[\textbf{Restituisce}]
    \texttt{pandas.DataFrame} con i dati caricati e validati.
  \item[\textbf{Eccezioni}]
    \texttt{ValueError} se \texttt{missing} non è una policy supportata;
    \texttt{ValueError} se una o più colonne richieste sono assenti;
    \texttt{ValueError} se il file contiene valori mancanti e \texttt{missing="error"}.
\end{apidesc}

\paragraph{Policy per i valori mancanti}
\begin{itemize}
  \item \texttt{missing="error"}: il caricamento fallisce se compare almeno un \texttt{NaN}
  \item \texttt{missing="drop"}: le righe incomplete vengono eliminate con \texttt{DataFrame.dropna()}
  \item \texttt{missing="allow"}: il \texttt{DataFrame} viene restituito così com'è
\end{itemize}

\paragraph{Implementazione essenziale}
\begin{minted}{python}
df = pd.read_csv(
    Path(path),
    sep=sep,
    decimal=decimal,
    comment=comment,
    skipinitialspace=skip_initial_space,
)

if rename_columns:
    df = df.rename(columns=rename_columns)

if required_columns:
    missing_columns = [col for col in required_columns if col not in df.columns]
    if missing_columns:
        raise ValueError(f"Colonne mancanti: {missing_columns}")
\end{minted}

\paragraph{Esempio d'uso}
\begin{minted}{python}
df = load_csv(
    "misure.csv",
    sep=";",
    decimal=",",
    rename_columns={"t [s]": "t", "x [m]": "x"},
    required_columns=["t", "x"],
    missing="drop",
)
\end{minted}
